% ============================================================
%  IEEE Conference Paper — IEEEtran template
%  Compile: pdflatex paper -> bibtex paper -> pdflatex paper (x2)
% ============================================================
\documentclass[conference]{IEEEtran}
\IEEEoverridecommandlockouts

% ---- packages -----------------------------------------------
\usepackage{cite}
\usepackage{amsmath,amssymb,amsfonts}
\usepackage{graphicx}
\usepackage{textcomp}
\usepackage{xcolor}
\usepackage{booktabs}
\usepackage{multirow}
\usepackage{array}
\usepackage{balance}       % balances columns on last page

\def\BibTeX{{\rm B\kern-.05em{\sc i\kern-.025em b}\kern-.08em
    T\kern-.1667em\lower.7ex\hbox{E}\kern-.125emX}}

% ============================================================
\begin{document}

\title{Reinforcement Learning for Adaptive Traffic Signal Control
at an Isolated Four-Way Intersection:\\
A Comparative Study of Tabular and Deep RL Algorithms}

\author{%
\IEEEauthorblockN{[Authors]}
\IEEEauthorblockA{%
\textit{[Department, Institution]}\\
[City, Country]\\
[email@institution.edu]}}

\maketitle

% ============================================================
% ABSTRACT
% ============================================================
\begin{abstract}
Urban traffic congestion remains a persistent challenge in modern
cities, and adaptive traffic signal control~(ATSC) via reinforcement
learning~(RL) has emerged as a promising solution. This paper
investigates RL-based ATSC at an isolated four-way intersection
under two action-space configurations: a simplified two-action setup
(North-South green versus East-West green) and a fine-grained
four-action movement-based setup that separately controls
straight-through and protected left-turn/U-turn phases for each
axis. Three classical tabular algorithms---TD(0), SARSA, and
Q-Learning---are benchmarked alongside two deep RL methods:
Pri-DDQN (priority-replay Double DQN)~\cite{Zheng2024PriDDQNLAG}
and a single-intersection adaptation of MA-SAC with a
spatio-temporal hypergraph-inspired centralised
critic~\cite{Wang2024TowardMRP}. All experiments are conducted in a
lightweight stochastic discrete-time simulator that captures Poisson
vehicle arrivals, phase-change yellow losses, and per-movement
service rates. Results demonstrate that deep RL methods achieve
substantially lower average queue lengths and waiting times than
tabular baselines, with Pri-DDQN attaining the best cumulative
reward in the four-action setting while both deep RL approaches
achieve near-identical traffic efficiency.
\end{abstract}

\begin{IEEEkeywords}
adaptive traffic signal control, reinforcement learning,
deep Q-network, soft actor-critic, Q-learning, SARSA,
four-way intersection, action space design
\end{IEEEkeywords}

% ============================================================
% 1. INTRODUCTION
% ============================================================
\section{Introduction}

Signalised intersections are the primary bottleneck in urban road
networks. Fixed-time controllers, while straightforward to deploy,
cannot adapt to the natural variability of traffic demand, leading
to unnecessary delays and fuel consumption. Adaptive traffic signal
control~(ATSC) has therefore been an active research area for
several decades, and reinforcement learning~(RL) has emerged as a
particularly attractive framework because it allows a controller to
learn an optimal phase-switching policy directly from interaction
with the environment, without requiring an explicit traffic model.

Early work established the viability of Q-learning for isolated
intersections~\cite{Abdulhai2003ReinforcementLFAD,
ElTantawy2010AnALV, A2011ReinforcementLWU}.
Abdulhai~et~al.~\cite{Abdulhai2003ReinforcementLFAD} demonstrated
that a single-agent Q-learning controller can outperform actuated
control at a simulated isolated intersection, and subsequent work
began to formalise state representation, action space, and reward
function as critical design
decisions~\cite{ElTantawy2014DesignORH, Touhbi2017AdaptiveTSM}.
El-Tantawy~et~al.~\cite{ElTantawy2014DesignORH} systematically
investigated these parameters and achieved up to 48\% savings in
average vehicle delay compared to pre-timed control.

With the rise of deep neural networks, deep RL methods quickly
supplanted tabular approaches for high-dimensional state
spaces~\cite{Haydari2020DeepRLC, Rasheed2020DeepRLBL}.
Haydari and Yilmaz~\cite{Haydari2020DeepRLC} provide a
comprehensive survey of deep RL applications in intelligent
transportation covering state encodings, reward formulations, and
simulation environments. Gao~et~al.~\cite{Gao2017AdaptiveTSA}
applied DQN with experience replay and a target network to a
four-way intersection, establishing a widely reproduced baseline.
Wan and Hwang~\cite{Wan2018ValuebasedDRE} proposed a value-based
deep Q-learning controller with a dynamic discount factor for
isolated intersection control, achieving a 20\% reduction in system
delay relative to fixed timing. Li~et~al.~\cite{Li2020AdaptiveTSN}
extended DQN to both single- and multi-intersection settings,
measuring performance via average waiting time.

Double DQN reduces the overestimation bias of vanilla DQN and has
been adopted widely for traffic signal
control~\cite{Gu2020DoubleDQD, Zheng2024PriDDQNLAG}.
Gu~et~al.~\cite{Gu2020DoubleDQD} introduced a dual-agent DDQN
architecture for a four-phase isolated intersection, demonstrating
improved traffic capacity in SUMO simulations.
Zheng~et~al.~\cite{Zheng2024PriDDQNLAG} further proposed Pri-DDQN,
a hybrid agent that integrates priority-based dynamic replay,
power-function epsilon decay, and state/reward-aware loss
weighting, reducing both average queue length and waiting time
at single intersections. Liao~et~al.~\cite{Liao2020TimeDPQ}
explored an LSTM Q-network with a time-difference penalty to
balance safety and throughput, reporting improvements over fixed
signal timing and vanilla DQN across eight signal phases.

Beyond standard DQN variants, hierarchical and policy-gradient
methods have been explored. Gu~et~al.~\cite{Gu2025AHDL} proposed a
hierarchical DDPG framework that separately allocates cycle
durations at NS and EW axes and then divides green time between
straight and protected left-turn movements---a structure closely
related to our four-action formulation.
Ma~et~al.~\cite{Ma2021AdaptiveOOK} applied Proximal Policy
Optimisation~(PPO) with Discrete Traffic State Encoding to
dynamically adjust phase lengths, and Pan~\cite{Pan2023TrafficLCR}
combined a phase-gate component with a memory palace mechanism to
improve queue-length and travel-time performance. Sekar and
Nezamoddini~\cite{Sekar2025OptimizingMIB} employed Soft
Actor-Critic~(SAC) for four-way intersection control in a mixed
autonomous/human-driven vehicle environment.

Multi-agent extensions have attracted considerable attention for
coordinating signals across intersection
networks~\cite{Balaji2010UrbanTSJ, ElTantawy2012MultiAgentRLY,
ElTantawy2013MultiagentRLAG, Wiering2000MultiAgentRLAF,
Xu2013TheSOAL, Kim2019CooperativeTSAB,
Kolat2023MultiAgentRLAN}.
El-Tantawy and Abdulhai~\cite{ElTantawy2012MultiAgentRLY} developed
MARLIN-ATSC for a five-intersection network, later scaled to
downtown Toronto with a 27\% delay
reduction~\cite{ElTantawy2013MultiagentRLAG}.
Wang~et~al.~\cite{Wang2024TowardMRP} introduced MA-SAC with a
spatio-temporal hypergraph critic capturing dependencies across
multiple intersections simultaneously---a design we adapt to the
single-intersection benchmark in this paper. BCT-APLight
\cite{2024BayesianCRT} further advanced multi-intersection ATSC
with an attention-based adaptive pressure mechanism.

Reward function design has received dedicated
attention~\cite{Touhbi2017AdaptiveTSM, Egea2020AssessmentORW,
Aslani2018ContinuousRRZ,Lee2022EffectsAOBK}. State and action
representation issues are discussed in Kanis~et~al.
\cite{Kanis2021BackTBO}, who show that acyclic phase transitions
outperform cyclic ones, and in Varela~et~al.
\cite{Varela2021AMFX}, who propose a standardised RL-ATSC
development methodology. Broader perspectives include
connected-vehicle safety~\cite{Karbasi2024ExploringTIF},
visible-light-communication
integration~\cite{Vieira2024EnhancingUIAC},
digital-twin frameworks~\cite{Bagabaldo2025DigitalTFS},
human-centric objectives~\cite{Vlachogiannis2023HumanLightIRI},
universal V2X frameworks~\cite{Wang2023UniTSAAUAE}, and
Transformer-based approaches~\cite{Zhao2024SequenceDTAH}.

In this paper we make the following contributions: (i)~we design a
lightweight stochastic RL environment with two distinct
action-space configurations---a coarse two-action setup and a
fine-grained four-action movement-based setup; (ii)~we conduct a
systematic comparison of three tabular RL algorithms
(TD(0), SARSA, Q-Learning) against two deep RL benchmarks
(Pri-DDQN and MA-SAC adapted); (iii)~we analyse the impact of
action-space granularity on learning efficiency, throughput, queue
length, and waiting time; and (iv)~we discuss practical insights
for choosing between tabular and deep RL controllers at real-world
isolated four-way intersections.

% ============================================================
% 2. RELATED WORK
% ============================================================
\section{Related Work}

\subsection{Tabular RL at Isolated Intersections}

The earliest RL applications to ATSC used Q-learning at simulated
isolated intersections~\cite{Abdulhai2003ReinforcementLFAD}.
El-Tantawy and Abdulhai~\cite{ElTantawy2010AnALV} extended this to
multiple state models under Q-learning with variable phasing.
Systematic investigation of RL design parameters---state space,
action space, reward, phasing scheme---was carried out by
El-Tantawy~et~al.~\cite{ElTantawy2014DesignORH}, who achieved 48\%
average delay savings. Touhbi~et~al.~\cite{Touhbi2017AdaptiveTSM}
analysed generalised state spaces and reward definitions using
Q-learning, demonstrating consistent performance gains over
fixed-plan control. Bhatnagar and
colleagues~\cite{A2011ReinforcementLWU} compared Q-learning with
policy gradient actor-critic methods for adaptive traffic light
control. Aslani~et~al.~\cite{Aslani2017AdaptiveTSAO} further
applied actor-critic methods with real-world traffic disruption
events, and Mannion~et~al.~\cite{Mannion2016AnERAP} provided an
experimental review covering multiple tabular and function-approximation
RL algorithms for ATSC.

\subsection{Deep RL for Single-Intersection Control}

Gao~et~al.~\cite{Gao2017AdaptiveTSA} applied DQN with experience
replay and a target network to a four-way intersection, targeting
vehicle delay reduction. Wan and Hwang~\cite{Wan2018ValuebasedDRE}
employed deep Q-learning with a dynamic discount factor, reducing
system delay by 20\%. Genders and Razavi~\cite{Genders2016UsingADAY}
were among the earliest to apply DNN-based DQN to traffic signals,
and later proposed asynchronous $n$-step
Q-learning~\cite{Genders2019AsynchronousNQAV} for improved
sample efficiency. Li~et~al.~\cite{Li2020AdaptiveTSN} built a
DQN-based ATSC model tested under varying traffic conditions.
Gu~et~al.~\cite{Gu2020DoubleDQD} introduced a dual-agent DDQN
with priority experience replay at a four-phase isolated
intersection. Zheng~et~al.~\cite{Zheng2024PriDDQNLAG} advanced
this further with Pri-DDQN, combining Double DQN with priority
replay and state/reward-aware loss weighting---a key baseline in
our study. Liao~et~al.~\cite{Liao2020TimeDPQ} proposed an LSTM Q-network
with a time-difference penalty. Pan~\cite{Pan2023TrafficLCR}
combined a phase-gate mechanism with a memory palace for DQL.
Jiang~et~al.~\cite{Jiang2022EconomicDrivenATAI} applied DDDQN with
an economic-driven penalty at four-way intersections.
Kodama~et~al.~\cite{Kodama2022TrafficSCAA} introduced a dual
targeting algorithm~(DTA) that reinforces successful experiences,
reporting a 33\% waiting-time reduction. Aslani~et~al.
\cite{Aslani2018ContinuousRRZ} used continuous residual RL with
feature functions to achieve a 15\% travel-time reduction.
Park~et~al.~\cite{Park2021DeepQTAM} provided an in-depth study of
DQN model variants for traffic signal control. Agafonov and
Myasnikov~\cite{Agafonov2021TrafficSCBQ} and Zeng~et~al.
\cite{Zeng2018AdaptiveTSAU} explored Double Q-learning and deep
recurrent Q-learning for adaptive signal timing, respectively.
Rodrigues and Azevedo~\cite{Rodrigues2019TowardsRDSJ} analysed
robustness of deep RL controllers under demand surges and sensor
failures. Survey and review articles covering DRL-based ATSC
include~\cite{Haydari2020DeepRLC, Rasheed2020DeepRLBL,
Yau2017ASOBO, Greguri2020ApplicationODAR, Wang2022ACRBC}.

\subsection{Hierarchical and Policy-Gradient Methods}

Ma~et~al.~\cite{Ma2021AdaptiveOOK} applied PPO with Discrete
Traffic State Encoding (DTSE) to dynamically adjust phase
durations. Gu~et~al.~\cite{Gu2025AHDL} proposed a two-level DDPG
framework for four-way intersections where the high-level agent
allocates time between NS and EW and the low-level agent
distributes green between straight and left-turn movements.
Mannion~et~al.~\cite{Mannion2015ParallelRLAK} used parallel RL to
increase exploration in complex multi-phase intersections.
Sekar and Nezamoddini~\cite{Sekar2025OptimizingMIB} applied SAC in
a mixed-autonomy four-way scene with safety constraints.
Zhao~et~al.~\cite{Zhao2024SequenceDTAH} proposed the Sequence
Decision Transformer (SDT) for offline decision-based ATSC.
Tan~et~al.~\cite{Tan2019DeepRLBH, Tan2022DeepRLBJ} applied deep RL
with actor-critic structures for adaptive signal timing.
Lin and Lin~\cite{Lin2025SmartCTAZ} combined STGCN-LSTM with PPO
for smart-city traffic optimisation. Bouktif~et~al.
\cite{Bouktif2021TrafficSCBM} proposed a hybrid discrete-continuous
action space for traffic signal control. Wang~et~al.
\cite{Wang2025APHBB} designed a parallel hybrid action space RL
model for real-world ATSC.

\subsection{Multi-Agent and Network-Level Control}

Balaji~et~al.~\cite{Balaji2010UrbanTSJ} deployed multi-agent
Q-learning across 29 urban intersections.
Wiering~\cite{Wiering2000MultiAgentRLAF} introduced multiple MARL
algorithms for coordinated traffic light control.
El-Tantawy and Abdulhai~\cite{ElTantawy2012MultiAgentRLY} developed
MARLIN-ATSC for five-intersection networks, later scaled to
downtown Toronto~\cite{ElTantawy2013MultiagentRLAG}.
Xu~et~al.~\cite{Xu2013TheSOAL} modelled multi-agent traffic control
as a Markov game using joint-action Q-learning.
Kim and Jeong~\cite{Kim2019CooperativeTSAB} proposed cooperative
DQN with traffic flow prediction. BCT-APLight~\cite{2024BayesianCRT}
applied Bayesian critique-tuning with adaptive pressure.
Wang~et~al.~\cite{Wang2024TowardMRP} proposed MA-SAC with a
spatio-temporal hypergraph critic; we adapt this architecture to
the single-intersection benchmark. Arel~et~al.
\cite{Arel2010ReinforcementLMBT}, Zhang~et~al.
\cite{Zhang2021IndependentRLAW}, Kolat~et~al.
\cite{Kolat2023MultiAgentRLAN}, Hu and Li~\cite{Hu2024AMDBA}, and
Jiang~et~al.~\cite{Jiang2022MultiAgentRLBV} present further MARL
extensions for network-level coordination.

% ============================================================
% 3. TRAFFIC ENVIRONMENT SETUP
% ============================================================
\section{Traffic Environment Setup}

All algorithms are evaluated in a lightweight stochastic
discrete-time simulator that abstracts the key dynamics of a single
4-way intersection without requiring a full microscopic simulator
such as SUMO. Two experimental configurations are studied.

\subsection{Common Intersection Geometry}

Both experiments model a \textbf{single isolated four-way
(four-legged) signalised intersection} with four arms: North~(N),
South~(S), East~(E), and West~(W). A single RL controller decides
which movement group receives green at every decision step.
Table~\ref{tab:common} summarises shared simulation parameters.

\begin{table}[h]
\centering
\caption{Shared Simulation Parameters}
\label{tab:common}
\begin{tabular}{ll}
\toprule
Parameter & Value \\
\midrule
Simulation horizon & 3{,}600\,s (1\,h) \\
Decision interval $\Delta t$ & 10\,s \\
Yellow clearance $\tau_y$ & 3\,s \\
Effective green (same phase) & 10\,s \\
Effective green (phase change) & $\max(10-3,1) = 7$\,s \\
Max queue per lane/group & 30\,veh \\
Arrival model & Poisson($\lambda_i \cdot \Delta t$) \\
Traffic demand mode & \texttt{normal} \\
Training episodes & 500 \\
Evaluation episodes & 30 (seeds 142--171) \\
\bottomrule
\end{tabular}
\end{table}

\subsection{Experiment~A: Two-Action Setup}

The intersection is simplified to four directional queues (one per
arm: N, S, E, W) with only two discrete signal phases:

\begin{center}
\begin{tabular}{cl}
\toprule
Action & Phase \\
\midrule
0 & \textbf{NS\_GREEN} -- N and S arms receive green \\
1 & \textbf{EW\_GREEN} -- E and W arms receive green \\
\bottomrule
\end{tabular}
\end{center}

\noindent\textbf{State space (9-dimensional):}
\begin{equation}
  s = \bigl[q_N,\,q_S,\,q_E,\,q_W,\;
              w_N,\,w_S,\,w_E,\,w_W,\;
              \phi\bigr]
\end{equation}
where $q_i = Q_i/Q_{\max}$, $w_i = \min(W_i/W_{\max}, 1)$,
$W_{\max}=300$\,s, and $\phi\in\{0,1\}$ is the current phase.%
\footnote{For tabular agents the state is further discretised into
5 bins per queue dimension plus the current phase (tuple length 5).}

\noindent\textbf{Arrival rates} (normal mode):
$\lambda_{\{N,S,E,W\}} = 0.4$\,veh/step.
Service: $\approx 0.5$\,veh/step Poisson departure per green lane.

\noindent\textbf{Reward} (clipped to $[-1,1]$):
\begin{equation}
  r = -0.50\,P_q -0.30\,P_w +0.10\,B_{tp}
      -0.05\,P_{sw} -0.05\,P_{fair}
\end{equation}
where $P_q{=}\sum_i Q_i/(4Q_{\max})$,
$P_w{=}\sum_i W_i/(4W_{\max})$,
$B_{tp}$ is the normalised throughput bonus,
$P_{sw}{=}0.1$ if phase changed (else 0), and
$P_{fair}{=}\sigma(Q)/Q_{\max}$.

\subsection{Experiment~B: Four-Action Movement-Based Setup}

Experiment B introduces eight movement groups (straight+right and
left+U-turn per arm) and four discrete actions, each granting green
to one protected stage:

\begin{table}[h]
\centering
\caption{Four-Action Movement Definitions}
\label{tab:actions}
\begin{tabular}{clcc}
\toprule
Action & Name & Served Groups & SUMO Phase \\
\midrule
0 & NS Straight    & n\_through, s\_through & 0 \\
1 & NS Left/U-turn & n\_turn, s\_turn       & 2 \\
2 & EW Straight    & e\_through, w\_through & 4 \\
3 & EW Left/U-turn & e\_turn, w\_turn       & 6 \\
\bottomrule
\end{tabular}
\end{table}

\noindent\textbf{State space (17-dimensional):}
\begin{equation}
  s = \bigl[q_1,\ldots,q_8,\;w_1,\ldots,w_8,\;\phi/3\bigr]
\end{equation}
following the order \{n\_through, n\_turn, s\_through, s\_turn,
e\_through, e\_turn, w\_through, w\_turn\}.

\noindent\textbf{Arrival rates} (normal mode):
$\lambda_{through}{=}0.10$, $\lambda_{turn}{=}0.02$\,veh/s per arm.
Service: $\mu_{through}{=}0.80$, $\mu_{turn}{=}0.55$\,veh/s.

\noindent\textbf{Reward} (clipped to $[-1,1]$):
\begin{equation}
  r = -0.45\,P_q -0.30\,P_w +0.20\,B_{tp}
      -0.05\,P_{fair} -0.05\,P_{sw}
\end{equation}
where terms are averaged over all eight movement groups,
$B_{tp}{=}N_{dep}/(2Q_{\max})$, and
$P_{sw}{=}\mathbf{1}[\phi \neq \phi_{prev}]$.

% ============================================================
% 4. METHODOLOGY
% ============================================================
\section{Methodology}

% [To be completed in future revision.]

% ============================================================
% 5. DISCUSSION AND RESULTS
% ============================================================
\section{Discussion and Results}

All results are from \textbf{greedy evaluation} (no exploration)
over 30~independent seeds under the \texttt{normal} traffic demand
mode. Training uses seed~42. We report mean~$\pm$~std where
available.

\subsection{Experiment~A: Two-Action Single 4-Way Intersection}

Table~\ref{tab:2action} reports evaluation results for the two-action
setup. Only deep RL methods are evaluated in this configuration.

\begin{table}[h]
\centering
\caption{Exp.~A: Two-Action Results (30 seeds, normal mode)}
\label{tab:2action}
\setlength{\tabcolsep}{4pt}
\begin{tabular}{lcccc}
\toprule
\textbf{Method} & \textbf{Reward} & \textbf{Avg Q} & \textbf{Avg Wait} & \textbf{Thput}\\
 & (mean$\pm$std) & (veh) & (s) & (veh/ep) \\
\midrule
Pri-DDQN~\cite{Zheng2024PriDDQNLAG}
  & $-58.33 \pm 2.30$ & 0.45 & 3.92 & 5{,}767 \\
MA-SAC (adapted)~\cite{Wang2024TowardMRP}
  & $-65.51 \pm 2.53$ & 1.00 & 4.75 & 5{,}762 \\
\bottomrule
\end{tabular}
\end{table}

In the two-action setting, Pri-DDQN achieves lower average queue
(0.45 vs.\ 1.00\,veh), lower waiting time (3.92 vs.\ 4.75\,s), and
a better mean reward ($-58.33$ vs.\ $-65.51$). Both methods
deliver virtually identical throughput ($\approx$5{,}762--5{,}767\,veh
per episode). The high throughput compared to Experiment~B reflects
the coarser action space: each two-action decision serves two full
arms simultaneously, accommodating higher aggregate flow per step.

\subsection{Experiment~B: Four-Action Movement-Based Intersection}

Table~\ref{tab:4action} presents all five methods ordered from
worst to best cumulative reward.

\begin{table}[h]
\centering
\caption{Exp.~B: Four-Action Results (30 seeds, normal mode)}
\label{tab:4action}
\setlength{\tabcolsep}{3pt}
\begin{tabular}{lccccc}
\toprule
\textbf{Method} & \textbf{Reward} & \textbf{Avg Q} & \textbf{Avg Wait} & \textbf{Thput} & \textbf{Ph.Ch.}\\
 & (mean$\pm$std) & (veh) & (s) & (veh/ep) & \\
\midrule
TD(0)
  & $-346.25\pm1.31$ & 19.00 & 1{,}359.5 & 718   & 0   \\
Q-Learning
  & $-216.21\pm73.83$ & 8.22 & 358.3    & 1{,}386 & 136 \\
SARSA
  & $-87.24\pm71.91$  & 3.30 & 103.9    & 1{,}634 & 218 \\
MA-SAC~\cite{Wang2024TowardMRP}
  & $-25.89\pm1.25$   & 1.08 & 21.6     & 1{,}722 & 330 \\
Pri-DDQN~\cite{Zheng2024PriDDQNLAG}
  & $-19.59\pm0.62$   & 0.97 & 21.3     & 1{,}717 & 244 \\
\bottomrule
\end{tabular}
\end{table}

\subsubsection{Tabular RL Analysis}
TD(0) converged to a degenerate policy (phase changes\,$=$\,0),
always selecting Action~0 (NS Straight). This produces catastrophic
East-West queue build-up: average queue 19.00\,veh, waiting time
1{,}359.5\,s, throughput just 718\,veh/episode. The failure stems
from TD(0) learning only $V(s)$, which cannot distinguish which
\emph{action} causes a transition, causing the greedy policy to
collapse to one fixed action. Q-Learning shows partial recovery but
high variance ($\pm73.83$), indicating convergence to different
local optima across seeds. SARSA outperforms Q-Learning
significantly---average queue 3.30\,veh, waiting time 103.9\,s,
throughput 1{,}634\,veh---because on-policy updates prevent
over-optimistic estimates for unvisited state-action pairs.

\subsubsection{Deep RL Analysis}
Pri-DDQN achieves the best reward ($-19.59$) and the lowest average
queue (0.97\,veh), a $3.4\times$ improvement over SARSA.
MA-SAC (adapted) delivers comparable traffic efficiency
(queue~1.08\,veh, wait~21.6\,s, throughput~1{,}722\,veh) with a
higher phase-change count (330 vs.\ 244), reflecting the continuous
actor's more dynamic switching behaviour. The reward gap between
Pri-DDQN and MA-SAC ($-19.59$ vs.\ $-25.89$) is at least partly
attributable to the additional switch penalty accumulated by MA-SAC.

\subsubsection{Key Observations}
\begin{enumerate}
  \item \textbf{TD(0) collapse:} A value-only method that mirrors
    $V(s)$ into action slots fails catastrophically in a 4-action
    environment and should not be used as a standalone ATSC
    controller.
  \item \textbf{SARSA vs.\ Q-Learning:} On-policy SARSA is
    substantially more robust than off-policy Q-Learning,
    achieving lower variance and better efficiency.
  \item \textbf{Tabular vs.\ deep RL gap:} The best tabular method
    (SARSA) is outperformed by deep RL on every metric:
    $3.4\times$ fewer queued vehicles, $4.9\times$ lower waiting
    time, and higher reward.
  \item \textbf{Action-space granularity:} The four-action
    movement-based formulation enables deep RL agents to reduce the
    average queue below one vehicle under normal demand, a
    resolution not achievable with the coarser two-action setup.
  \item \textbf{Pri-DDQN vs.\ MA-SAC:} Both achieve near-identical
    traffic efficiency; Pri-DDQN slightly outperforms on reward
    owing to fewer unnecessary phase changes.
\end{enumerate}

% ============================================================
% 6. CONCLUSION AND FUTURE WORK
% ============================================================
\section{Conclusion and Future Work}

% [To be completed in future revision.]

% ============================================================
% REFERENCES — use external papers.bib
% ============================================================
\balance
\bibliographystyle{IEEEtran}
\bibliography{papers}

\end{document}
